% Generated draft: requires mathematical and source review.
\documentclass[11pt]{article}
\usepackage[margin=25mm]{geometry}
\usepackage{fontspec}
\setmainfont{Latin Modern Roman}
\usepackage{amsmath,amssymb,mathtools,unicode-math}
\setmathfont{Latin Modern Math}
\usepackage{enumitem,needspace,xurl,hyperref}
\setlength{\parindent}{0pt}
\setlength{\parskip}{5pt}
\setlength{\emergencystretch}{3em}
\newcommand{\N}{\mathbb{N}}
\newcommand{\Z}{\mathbb{Z}}
\newcommand{\Q}{\mathbb{Q}}
\newcommand{\R}{\mathbb{R}}
\newcommand{\C}{\mathbb{C}}
\newcommand{\F}{\mathbb{F}}
\newcommand{\A}{\mathbb{A}}
\newcommand{\PP}{\mathbb{P}}
\newcommand{\E}{\mathbb{E}}
\newcommand{\eps}{\varepsilon}
\newcommand{\dd}{\,\mathrm{d}}
\begin{document}
\section*{UnsolvedMath ID: 30001058}
\textbf{Source problem code:} catalogue entry
\subsection{Definitions and mathematical statement}
Fix $d\geq1$. For every $h(t)=\sum_{i=0}^{d}h_it^i$ with $h_i\in\Z_{\geq0}$ and $h_0\geq1$, define $g(m)$ and the Veronese numerator $U_nh$ by
\[
\frac{h(t)}{(1-t)^{d+1}}=\sum_{m\geq0}g(m)t^m,
\qquad
\frac{U_nh(t)}{(1-t)^{d+1}}=\sum_{m\geq0}g(nm)t^m.
\]
Write $U_nh(t)=\sum_{i=0}^{d}h_i(n)t^i$, allowing the original $h$ to have degree below $d$.

Determine optimal dimension-dependent positive integers $m_d,n_d$ for which, for every such $h$ and every integer $n\geq n_d$, all of the following hold:
\begin{enumerate}
\item $U_nh$ has $d$ negative real roots, and for each fixed $h$ their multiset converges as $n\to\infty$ to that of the intended Eulerian polynomial $A_d(t)$;
\item all coefficients $h_i(n)$ are positive, and $h_i(n)^2>h_{i-1}(n)h_{i+1}(n)$ for $1\leq i\leq d-1$;
\item $h_i(n)<m_dh_d(n)$ for every $0\leq i\leq d$.
\end{enumerate}
A degree-$d$ Eulerian convention compatible with the root count is $A_d(t)=\sum_{j=1}^{d}\left\langle\begin{smallmatrix}d\\j-1\end{smallmatrix}\right\rangle t^j$, using Eulerian numbers; it includes a zero root as a possible limit of negative roots. The extract does not define its Eulerian normalization, so this convention requires source confirmation.

Retain also the Ehrhart special case. For a $d$-dimensional lattice polytope $P$, define $\delta_P(t)$ by
\[
\sum_{r\geq0}|rP\cap\Z^d|t^r=\frac{\delta_P(t)}{(1-t)^{d+1}}.
\]
Then $U_n\delta_P=\delta_{nP}$. Determine the optimal corresponding constants for the source's real-rootedness, positivity, log-concavity and asymptotic-root conclusions together with
\[
\delta_0(n)=1,\qquad
\delta_i(n)<\delta_{d-i}(n)<\delta_{i+1}(n)
\quad\left(0\leq i\leq\left\lfloor\frac d2\right\rfloor-1\right),
\]
\[
\delta_{\lfloor(d+1)/2\rfloor}(n)<m_d\delta_d(n).
\]
The constants must be uniform over the indicated input class. The source does not define whether ``optimal'' minimizes one constant with the other fixed, minimizes both simultaneously, or describes the tradeoff of admissible pairs, nor whether the restricted Ehrhart class uses the same optimal pair as the general class. Those optimization conventions are not guessed.
\subsection{Short English statement}
Find sharp dimension-only thresholds and coefficient constants for Veronese numerators and Ehrhart dilation polynomials, retaining all real-root, log-concavity and interlacing inequalities. The Eulerian normalization and the precise optimization of the two constants need clarification.
\subsection{Sources}
\begin{itemize}
\item Supplied report: \emph{Discrete Geometry} (2009), Oberwolfach Reports, DOI 10.4171/owr/2008/44, pp. 2477--2556; proposed year 2008. \url{https://doi.org/10.4171/owr/2008/44}.
\item Supplied Beck--Stapledon reference: \url{https://arxiv.org/abs/0804.3639}.
\item Local UnsolvedMath registry, permanent ID 30001058, catalogue entry; created 25 August 2026, updated 16 September 2026. Its corrected\_verified note reports preservation of the theorem/corollary hypotheses and optimization target after checking the official report. Literature checked at \texttt{2026-08-21}: the assessment is supplied dated curation, not independently verified. It labels optimality open while reporting existence of dimension-dependent constants, distinguishing that from determining sharp values.
\end{itemize}
\end{document}
